\chapter{\conclusionsname}
\label{chap:conclusiones}

\section{Contribuciones del trabajo}

Este trabajo ha realizado un mapeo sistem\'atico y una evaluaci\'on emp\'irica del flujo de trabajo de \textsc{SpecKit} contra un flujo profesional de ingenier\'ia de requisitos basado en normas IREB, ISO 29148 e INCOSE. Las contribuciones principales son:

\begin{enumerate}
    \item \textbf{IREB Kit for SpecKit}: un conjunto de 10 artefactos (constitution con 9 art\'iculos, plantilla de spec con 12 atributos, checklist ISO 29148 de 13 criterios, AGENTS.md con 5 reglas universales, scope contract, y plantillas de plan, tasks, clarify y analyze) que transforman el pipeline nativo de \textsc{SpecKit} en un flujo alineado con pr\'acticas normativas de RE.

    \item \textbf{Corpus de 52 requisitos formalizados}: extra\'idos de \emph{changelogs} de 6 proyectos \emph{Open Source} representativos de distintos tipos de sistema, formalizados con 12 atributos IREB y validados mediante r\'ubrica de calidad. Complementado con 52 MRS (\emph{Minimal Requirement Specifications}) derivados para replicabilidad.

    \item \textbf{Infraestructura de evaluaci\'on autom\'atica}: scripts que aplican 6 herramientas (semgrep, eslint, phpcs, shellcheck, pyflakes y cloc) y comparan el c\'odigo generado con los PR reales mediante m\'etricas de precisi\'on, recall y F1. La r\'ubrica SRCI v3, dise\~nada como complemento cualitativo, queda preparada para su aplicaci\'on en la siguiente fase.

    \item \textbf{Evaluaci\'on emp\'irica de 4 flujos sobre 24 casos}: 96 ejecuciones que muestran c\'omo cada nivel de madurez metodol\'ogica mejora la alineaci\'on con el PR real: C0 (F1=0.46) $\rightarrow$ C1 (+IREB, F1=0.52) $\rightarrow$ C2 (+control de alcance, F1=0.62). C3, sin pipeline, alcanza el mismo F1 que C2 (0.62) pero con un perfil distinto: mucha precisi\'on (83\,\%) sobre pocos archivos (3 de media), frente al equilibrio de C2 (70\,\% de precisi\'on sobre 15 archivos).

    \item \textbf{Mapa de alineaci\'on IREB--SpecKit}: contraste detallado de cada actividad IREB (7), cada criterio ISO 29148 (8) y cada principio IREB (9) contra los comandos y artefactos de \textsc{SpecKit}, con y sin el IREB Kit.
\end{enumerate}

\section{Conclusiones principales}

\textbf{1. A\~nadir estructura al proceso de generaci\'on mejora los resultados de forma consistente.} La progresi\'on de F1, desde C0 (0.46) pasando por C1 (0.52) hasta C2 (0.62), muestra que cada nivel adicional de formalizaci\'on y control acerca el c\'odigo generado al cambio real de referencia. La precisi\'on pasa del 49.5\,\% al 70.2\,\%, y la preservaci\'on de atributos IREB en la especificaci\'on sube de 1.5 a 7.0 sobre 7. La mejora es acumulativa: cada componente del IREB Kit contribuye de forma aditiva.

\textbf{2. El pipeline determina el alcance de la implementaci\'on generada.} C3 y C2 comparten el mismo F1 medio (0.62), pero lo consiguen por caminos complementarios. C3 destaca en precisi\'on (83.0\,\%) con una media de 3 archivos generados; C2 ofrece un equilibrio entre precisi\'on (70.2\,\%) y cobertura (62.5\,\%) sobre 15 archivos. Ambos perfiles resultan \'utiles en contextos distintos: C3 para cambios de alcance reducido y autocontenidos, C2 para requisitos que requieren trazabilidad y una cobertura amplia del cambio.

\textbf{3. El control de alcance es el mecanismo con mayor impacto individual.} El salto C1 $\rightarrow$ C2 (+0.10 en F1, +17 puntos porcentuales en precisi\'on) supera al salto C0 $\rightarrow$ C1 (+0.06 en F1). El \emph{scope contract} y las reglas de \texttt{AGENTS.md} no solo mejoran la precisi\'on: tambi\'en reducen el coste (de \$0.0915 a \$0.0763) y acotan el c\'odigo generado (de 19.696 a 12.491 l\'ineas), porque el agente consume menos tokens en c\'odigo no relacionado con el requisito.

\textbf{4. La ingenier\'ia de requisitos y SDD se complementan.} \textsc{SpecKit} sin IREB Kit produce c\'odigo funcional y cubre 14 archivos de media. Al a\~nadir requisitos formalizados, la precisi\'on, la cobertura y la trazabilidad mejoran en todas las m\'etricas. El ingeniero de requisitos elicita y formaliza; \textsc{SpecKit} implementa con trazabilidad. La combinaci\'on de ambos enfoques ofrece mejores resultados que cualquiera de ellos por separado.

\textbf{5. Unos requisitos claros y bien definidos ayudan al agente a no desviarse del objetivo.} En el caso de Authentik, C0 implement\'o una funcionalidad distinta a la solicitada porque el \emph{changelog} no proporcionaba suficiente contexto. C2, al recibir un requisito formalizado con criterios de verificaci\'on y reglas de comportamiento, gener\'o exactamente lo esperado. Dedicar tiempo a escribir buenos requisitos reduce el riesgo de que el agente malinterprete la tarea.

\textbf{6. El impacto de la formalizaci\'on IREB depende de la calidad de la fuente original.} Los repositorios cuyos \emph{changelogs} son menos descriptivos, como Medusa (+0.39 en F1) y Appwrite (+0.34), obtienen el mayor beneficio de la formalizaci\'on. En el extremo opuesto, Authentik ($-$0.09) ya alcanza un F1 elevado sin IREB porque sus \emph{changelogs} contienen descripciones t\'ecnicamente precisas. Esto indica que la formalizaci\'on IREB no es una pr\'actica de talla \'unica: su valor es directamente proporcional a la ambig\"uedad de la fuente de requisitos original.

\section{Resumen comparativo}

La tabla~6.5 del cap\'itulo anterior recoge todas las m\'etricas. En resumen: C2 lidera en F1 (0.62), precisi\'on (70.2\,\%) y preservaci\'on de atributos IREB (7.0/7); C1 ofrece el mejor recall (64.8\,\%); C3 es la opci\'on m\'as econ\'omica (\$0.0191) y de mayor precisi\'on (83.0\,\%), con un alcance m\'as reducido (3 archivos de media); C0 destaca por su bajo coste dentro de los flujos con pipeline (14 archivos a \$0.0573) y por generar el c\'odigo con menos incidencias de estilo (126). La elecci\'on depende del contexto: para cambios peque\~nos y autocontenidos, C3 resulta adecuado; para requisitos que necesitan trazabilidad y cobertura, C2 ofrece las mayores garant\'ias.

\section{Trabajo futuro}

\begin{enumerate}
    \item \textbf{Evaluaci\'on manual SRCI v3}: aplicar la r\'ubrica de inspecci\'on de requisitos a los 96 artefactos generados, que ya est\'an disponibles y no requieren repetir los experimentos.

    \item \textbf{M\'ultiples ejecuciones por caso}: repetir cada configuraci\'on 3--5 veces para medir la variabilidad del modelo y realizar an\'alisis estad\'istico de las diferencias observadas.

    \item \textbf{Ampliar el corpus de ejecuci\'on}: extender la evaluaci\'on a los 28 requisitos del corpus que no se han ejecutado, para comprobar si los resultados se mantienen en una muestra m\'as amplia.

    \item \textbf{Explorar otros modelos}: repetir los experimentos con modelos alternativos para verificar si las diferencias entre flujos son consistentes independientemente del modelo de lenguaje utilizado.

    \item \textbf{Integraci\'on con herramientas profesionales}: estudiar la viabilidad de conectar el IREB Kit con herramientas como DOORS o Jama para cerrar el ciclo completo desde la gesti\'on de requisitos hasta la generaci\'on de c\'odigo.
\end{enumerate}

\section{Reflexi\'on final}

Este trabajo comenz\'o con una pregunta concreta: ¿qu\'e ocurre al integrar la ingenier\'ia de requisitos con las herramientas de generaci\'on de c\'odigo asistida por IA? La respuesta, respaldada por 96 ejecuciones sobre 24 requisitos reales, es que ambas disciplinas se potencian mutuamente. Un pipeline sin requisitos formalizados (C0) ya ofrece una base s\'olida de partida. Formalizar esos requisitos (C1, C2) mejora todos los indicadores. Incluso sin pipeline (C3), el modelo es capaz de producir implementaciones precisas, aunque de alcance limitado. La conclusi\'on principal es que el camino m\'as productivo no consiste en elegir entre ingenier\'ia de requisitos y generaci\'on autom\'atica, sino en combinarlas: el ingeniero aporta el rigor en la definici\'on de qu\'e debe hacerse, y el pipeline estructurado se encarga de transformar esa definici\'on en c\'odigo con trazabilidad y precisi\'on.